%! Change in Linear and Exponential Functions — Unit 4, Lesson 4.2 — Guided Notes
\documentclass[10pt]{article}
\usepackage{precalculus-article}
\usepackage{precalculus-boxes}
\newcommand{\TallMath}[1]{$\displaystyle #1\rule[-1.4em]{0pt}{3.2em}$}

\begin{document}

\pageheader{Unit 4, Lesson 4.2}{Guided Notes}
\namedateperiod

\vspace{0.06in}

\begin{objectivebox}
By the end of this lesson, I will be able to\ldots
\begin{enumerate}[itemsep=2pt]
  \item Construct a \blank{3cm} function from a slope and a point using
        $f(x) = y_i + m(x - x_i)$. \hfill\textit{(LO 2.2.A)}
  \item Construct an \blank{3cm} function from a ratio and a point using
        $f(x) = y_i \cdot r^{x - x_i}$. \hfill\textit{(LO 2.2.A)}
  \item Determine whether outputs over equal-length intervals show a constant \blank{2.5cm}
        (linear) or a constant \blank{2cm} (exponential). \hfill\textit{(LO 2.2.B)}
\end{enumerate}
\end{objectivebox}

\vspace{0.04in}

\begin{vocabbox}
\termblanklong{linear function}
\termblanklong{exponential function}
\termblanklong{constant rate of change}
\termblanklong{constant ratio}
\termblanklong{initial value}
\termblanklong{domain}
\end{vocabbox}

\vspace{0.04in}

\begin{hookbox}
Two tables are presented without labels. For each table, analyze how the outputs change.

\medskip
\renewcommand{\arraystretch}{1.3}
\begin{tabular}{c|c}
$x$ & $f(x)$ \\ \hline
0 & 3 \\ 1 & 7 \\ 2 & 11 \\ 3 & 15
\end{tabular}
\qquad\qquad
\begin{tabular}{c|c}
$x$ & $g(x)$ \\ \hline
0 & 2 \\ 1 & 6 \\ 2 & 18 \\ 3 & 54
\end{tabular}

\medskip
Table 1: consecutive outputs differ by \blank{2cm}. \quad
Table 2: consecutive outputs are multiplied by \blank{2cm}.
\end{hookbox}

\vspace{0.04in}

\begin{notesbox}{LO 2.2.A --- Linear Functions and Arithmetic Sequences}
\renewcommand{\arraystretch}{1.5}
\begin{tabularx}{\linewidth}{@{} >{\bfseries}p{0.24\linewidth} X X @{}}
 & \textbf{Arithmetic Sequence} & \textbf{Linear Function} \\
\midrule
Formula & $a_n = a_0 + d\,n$ & $f(x) = b + mx$ \\[3pt]
From one known value & $a_n = a_i + d(n - i)$ & $f(x) = y_i + m(x - x_i)$ \\[3pt]
``Change'' quantity & common difference $d$ & slope $m$ \\[3pt]
Domain & \blank{3.5cm} & \blank{3.5cm} \\
\end{tabularx}

\medskip
\textbf{Example:} Slope $m = 4$; passes through $(1,\ 7)$.

\[
  f(x) = \blank{2cm} + \blank{2cm}(x - \blank{1.5cm})
       = \blank{6cm}
\]

Verify: $f(0) = $ \blank{2cm}, \quad $f(2) = $ \blank{2cm}.
\end{notesbox}

\vspace{0.04in}

\begin{notesbox}{LO 2.2.A --- Exponential Functions and Geometric Sequences}
\renewcommand{\arraystretch}{1.5}
\begin{tabularx}{\linewidth}{@{} >{\bfseries}p{0.24\linewidth} X X @{}}
 & \textbf{Geometric Sequence} & \textbf{Exponential Function} \\
\midrule
Formula & $g_n = g_0 \cdot r^n$ & $f(x) = a \cdot b^x$ \\[3pt]
From one known value & $g_n = g_i \cdot r^{n - i}$ & $f(x) = y_i \cdot r^{x - x_i}$ \\[3pt]
``Change'' quantity & common ratio $r$ & base $b$ \\[3pt]
Domain & \blank{3.5cm} & \blank{3.5cm} \\
\end{tabularx}

\medskip
\textbf{Example:} Ratio $r = 3$; passes through $(0,\ 2)$.

\[
  f(x) = \blank{2cm} \cdot \blank{2cm}^{x - \blank{1.5cm}}
       = \blank{6cm}
\]

Verify: $f(1) = $ \blank{2cm}, \quad $f(3) = $ \blank{2cm}.
\end{notesbox}

\vspace{0.04in}

\begin{notesbox}{LO 2.2.B --- The Equal-Length Interval Test}
Over equal-length input intervals $\Delta x$:

\medskip
\renewcommand{\arraystretch}{1.6}
\begin{tabularx}{\linewidth}{@{} >{\bfseries}p{0.22\linewidth} X X @{}}
 & \textbf{Linear} & \textbf{Exponential} \\
\midrule
Output pattern & differences \blank{3cm} & ratios \blank{3cm} \\[3pt]
How to check & subtract consecutive outputs & divide consecutive outputs \\[3pt]
Formula built from 2 pts & $m = \dfrac{y_2 - y_1}{x_2 - x_1}$ & $b = \left(\dfrac{y_2}{y_1}\right)^{1/(x_2 - x_1)}$ \\
\end{tabularx}

\medskip
\textbf{Key fact (EK 2.2.B.3):} Two distinct input-output pairs \blank{6cm} determine
either type of function.
\end{notesbox}

\end{document}
